\section{The compensated low estimate}\label{sec:compensated-low}

We now bound the finite compensated probe defined in
Equation~\eqref{eq:compensated-probe-definition}.  Its exact low separation
has two factors.  We first bound the completed row after summing the marked
slots by applying Lemma~\ref{lem:reflected-energy} directly.  We then prove
a quantitative Gram bound for the additive polynomial from
Section~\ref{sec:probe}.  Cauchy--Schwarz will combine these estimates to
give the exponent $3/16$.

\subsection{The marked completed row}

For a fixed rescaled subset $J$, with marked subset $J^c$, put
\begin{equation}
 \begin{gathered}
 d=\sum_{i\in J}\ell_i,\qquad
 X'=X/q_{p_J}\asymp XZ^{-d},\qquad
 Y'=Y/q_{p_J}\asymp YZ^{-d},\\
 M'=M-2d,\qquad \ell'=\ell-d,\qquad 0\le d\le\ell=1/6.
 \end{gathered}
 \label{old-eq:5.2}
\end{equation}
For a squarefree product $D$ of slot primes, define the marked completed
series
\[
 T_D(s_0,\psi)=\sum_{\substack{c\ {\rm sf}\\n}}
 \gamma_2(c)\overline{\alpha(cn^3)}\psi(cn^3)
 1_{D\mid cn^3}q_c^{-s_0}q_n^{-3s_0+1/2}.
\]
The ideals $c,n$ are the original ideals prime to $S$, represented by their
primary generators, with $c$ squarefree;
they may share primes.  All character powers retain their original zero
extensions.  For the empty product $D=1$ this is the completed $T(s_0,\psi)$.
For this fixed $J$, let
\[
\begin{split}
 B^J_{m,\sigma}(Z)=
 \sum_{p_i\in\mathcal P_i(Z),\ i\notin J}
 \prod_{i\notin J}\overline{\eta(p_i)}W_i(q_{p_i}/P_i)
 \sum_{\theta\in\widehat T}a_\theta\frac1{2\pi i}
 \int_{(4)} &(Zq_{p_{J^c}})^t\Phi(t)\\
 &\cdot T_{p_{J^c}}(t+1/2,
             \nu_\sigma\theta\chi_\bullet(m))\,dt.
\end{split}
\]
Thus $B^J$ is exactly the completed-row factor for the $J$ summand
of Equation~\eqref{eq:compensated-probe-definition} after the
marked slots have been summed.  The fixed tuple $p_J$ affects
only $X',Y'$ in its low separation.

\begin{lemma}[The compensated completed-row norm]\label{lem:probe-row-norm}
For every $\epsilon>0$, every fixed rescaled subset $J$, and every
$\sigma\in T$, with $Q=q_{b_*}X'Y'\asymp Z^{M'}$,
\begin{equation}
 \sum_{0<q_m\ll Q}|B^J_{m,\sigma}(Z)|^2
       \ll Z^{M'+((d-1/6)/4)_++\epsilon}
       =Z^{M'+\epsilon}.
 \label{old-eq:5.4}
\end{equation}
The bound is uniform in the fixed rescaled tuple and has no
loss depending on the mesh of the surviving slots.
\end{lemma}

\begin{proof}
Put $A=cn^3$ and $\varrho_i=q_{p_i}/P_i$ for $i\notin J$.  Since
$q_A=q_cq_n^3$ even when $c,n$ share primes, and
$q_{p_{J^c}}=Z^{\ell'}\prod_{i\notin J}\varrho_i$, Mellin inversion
of the displayed integral defining $B^J$ gives the physical smooth factor
\[
 \begin{gathered}
 V\!\left(\frac{q_A}{Zq_{p_{J^c}}}\right)
 =W_{\mathrm G}\!\left(\frac{r}{\prod_{i\notin J}\varrho_i}\right),\\
 r=\frac{q_A}{Z^{1+\ell'}}.
 \end{gathered}
\]
Here $V=W_{\mathrm G}$ is the Gaussian of
Lemma~\ref{lem:gaussian-annular}.  Empty products are one.

Use the fixed translate partition $\chi$ from that lemma on $r$.
On the annulus $r=e^kx$, retain the single joint profile
\[
 w_{k,J}(x,\boldsymbol\varrho)
 =\chi(\log x)W_{\mathrm G}\!\left(
       \frac{e^kx}{\prod_{i\notin J}\varrho_i}\right)
       \prod_{i\notin J}W_i(\varrho_i).
\]
The logarithms of $x$ and of all the $\varrho_i$ lie in fixed compact
sets on this support.  Thus the logarithm of the Gaussian argument is
$k+O_{\rm fixed}(1)$.  The derivative estimate in
Lemma~\ref{lem:gaussian-annular}, with this bounded shift and the fixed
slot windows, gives for every fixed $\kappa_{\mathrm G}>0$,
$A_{\rm cnt},B_{\rm cnt},C_{\rm ann}\ge0$, $D_{\rm tail}>0$, and
fixed seminorm order $j$,
\[
 \begin{gathered}
 \sum_{k\in\mathbb Z}e^{A_{\rm cnt}|k|}p_j(w_{k,J})<\infty,
 \\
 Z^{B_{\rm cnt}}
 \sum_{|k|>\kappa_{\mathrm G}\log Z-C_{\rm ann}}
       e^{A_{\rm cnt}|k|}p_j(w_{k,J})
       \ll Z^{-D_{\rm tail}}.
 \end{gathered}
\]
The constants may depend on the fixed slot system, but not on its moving
prime labels.  On an annulus $\log r=k+O(1)$, an absolute bound for the
row $\ell^2$ norm is at most
$Z^{B_{\rm cnt}}e^{A_{\rm cnt}|k|}p_j(w_{k,J})$ for some fixed exponents,
by the elementary row, slot, and ideal counts.  Choose $C_{\rm ann}$ larger
than the fixed logarithmic radius of $\chi$ and discard only whole annuli
with $|k|>\kappa_{\mathrm G}\log Z+C_{\rm ann}$.  Summing their norm bounds
and then squaring gives less than any prescribed power of $Z^{-1}$.
Every annulus meeting $|\log_Z r|\le\kappa_{\mathrm G}$ is retained.

For each remaining annulus, keep fixed individual annular cutoffs equal to
one on the support of $w_{k,J}$ and apply
Lemma~\ref{lem:smooth-calculus} to this whole profile.  Its logarithmic
Fourier inversion uses one coefficient density common to every row and
every surviving slot tuple.  The normalized slot ratios remain variables
of the joint profile until separation; they do not index separately
chosen densities.  For fixed Fourier variables, each resulting slot
coefficient is a product of the
original arithmetic coefficient $\overline{\eta(p_i)}$ and an individual
annular cutoff and norm power.  It remains bounded and independent of the
row and of every other slot.  The completed factor is an annular test at
length
\[
 N_*=1+\ell'+\theta_N,\qquad
 \theta_N=\frac{k}{\log Z},\qquad
 |\theta_N|\le\kappa_{\mathrm G}+O_{\rm fixed}(1/\log Z).
\]
These retained lengths lie in a fixed bounded range.  The weighted Fourier
norms have a summable total by the displayed Gaussian estimate; the small
annular weight stays in this homogeneous single-profile Fourier norm, not
in an inhomogeneous tuple seminorm.  Minkowski's inequality passes the
separated row norm through this common density.  The existing $q_c^{-1/2}q_n^{-1}$
normalization stays unchanged when the test scale is recentered, so no
additional power of $e^k$ is introduced.  Choose $\kappa_{\mathrm G}$
within the final $\epsilon$ allowance.  For each fixed $\theta$,
Lemma~\ref{lem:reflection-row-sectors} supplies the row-sector reduction of
$\nu_\sigma\theta\chi_\bullet(m)$ with every original zero mask.  It covers
all nonzero element rows, including those meeting $S$; no factor $\xi(m)$ is
used in this row norm.  The row ball depends only on the fixed rescaled tuple
$p_J$, through $Q=q_{b_*}X'Y'$, and is independent of the surviving marked
slots and dual variables.  The physical slot sets remain disjoint and their
separated coefficients are product-form.  Finite triangle over $\theta$ and
Lemma~\ref{lem:reflected-energy} therefore give the exponent in
Equation~\eqref{old-eq:4.6} at this actual $N_*$, taking the retention
parameter $\tau_{\rm ref}$ and the independently prescribed output
and separation losses all at most a small $\epsilon_0>0$.  We use
the exponent $E_{\rm ref}$ after the fixed row parts are summed;
their count is already included in $E_{\rm ref}+\epsilon_0$, with
principal exponent $O/2$, and is not counted again.

Here $f=1$ and the puncture is $\rho=1$, so the only moving non-slot
primes are those of the row.  After the finite unit and reflection-sector
splits in Lemma~\ref{lem:reflected-energy}, freeze its powerful part of norm length
$O$ and its valuation-one part supported on $S$, and split its residual
squarefree good product into \emph{actual} dyads of norm $\asymp Z^H$.  Let
$\epsilon_Z=O_{\rm fixed}(1/\log Z)$ be nonnegative and large enough
to contain all fixed annular and fixed-conductor logarithmic offsets
in this calculation.  Then
\[
 H\le M'-O+\epsilon_Z,\qquad
 \Delta_H=M'-O-H\ge-\epsilon_Z,\qquad O\ge-\epsilon_Z.
\]
The actual row and dual annular cutoffs are kept in one joint
profile while its common Fourier density is separated.  The small
kernel amplitude is extracted from that homogeneous density before
the row norm is squared, and only then may the positive sieve sum be
enlarged, as in Equation~\eqref{eq:reflection-common-minkowski}.
In particular the kernel saving in
Equation~\eqref{old-eq:4.6} uses this $H$; it is not reevaluated on
any shorter rows added by the positive enlargement.

In the notation of that energy estimate, every moving non-slot prime
counted in $A_0$ divides the powerful row part to exponent at least
two.  The comparison of its exact norm with the powerful dyad, and
any fixed-conductor contribution, therefore gives
\[
 2A_0\le O+\epsilon_Z,\qquad N_0\le A_0,\qquad z_a\le\ell'.
\]
The available dual length at the actual completed scale is
\[
 T_d=2H+2A_0+2z_a-1-\ell'-\theta_N-N_0-3B_0.
\]
Since $M'+\ell'-1=-3d$, direct subtraction gives
\[
 T_d-(H-3d)
 =H-M'+2A_0+2(z_a-\ell')-N_0-3B_0-\theta_N
 \le 2\epsilon_Z+|\theta_N|.
\]
As in the proof of the reflected energy, dual ranges below a fixed
negative length are negligible after a positive Mellin shift, and
the bounded boundary range costs an arbitrarily small power.
The source dual ideals $n,b$ in these dyads retain the support of
Equation~\eqref{old-eq:4.3}: they may share primes and may contain
the permitted primes of $S$; no additional $S$-mask or coprimality
condition between $n$ and $b$ is imposed.
On each remaining dual dyad, put $y=v+3\ell_b+e_\lambda$.  Retention
gives $y\le T_d+\epsilon_0$, while the nonzero unit ranges give
$v,\ell_b,e_\lambda\ge-\epsilon_Z$ after increasing its fixed
constant.  It follows that
\[
 \max(H,v+\ell_b)=H+O(\epsilon_0+\kappa_{\mathrm G}+\epsilon_Z).
\]

Write $s_{\rm hyb}=\min\{v,z_a,(v+z_a)/3\}$ for the quantity
called $u$ in Equation~\eqref{old-eq:4.7}.  If $s_{\rm hyb}=z_a$,
then $-S_0-B_0+z_a-s_{\rm hyb}\le0$, the kernel term is
nonpositive, and
\[
 -\ell_b-\frac{2e_\lambda}{3}\le\frac53\epsilon_Z.
\]
Here the standalone $O$ is the powerful-row logarithmic length.  It follows that
\[
 \begin{aligned}
 E_{\rm ref}&\le O/2+H+O(\epsilon_0+\kappa_{\mathrm G}+\epsilon_Z)\\
 &=M'-O/2-\Delta_H+O(\epsilon_0+\kappa_{\mathrm G}+\epsilon_Z).
 \end{aligned}
\]
Otherwise $s_{\rm hyb}\ge v/2-O(\epsilon_Z)$, and the bounded unit
ranges imply
\[
 s_{\rm hyb}+\ell_b+\frac{2e_\lambda}{3}
       +\frac{(T_d-y)_+}{2}
 \ge \frac y4+\frac{(T_d-y)_+}{2}-O(\epsilon_Z)
 \ge \frac{T_d}{4}-O(\epsilon_Z).
\]
The last inequality holds separately for $y\le T_d$ and $y\ge T_d$.
Here the squared kernel saving is obtained by first extracting
$\min\{1,Z^{(y-T_d)/4}\}$ from the common Fourier density and only
then squaring the row norm.  Relative to $M'$, the saving
before the remaining $+z_a$ is exactly
\begin{equation}
 \begin{split}
 O/2+\Delta_H+S_0+B_0+T_d/4
 ={}&\frac{2M'+2z_a-1-\ell'+2\Delta_H-\theta_N}{4}\\
    &+\frac{2A_0-N_0+B_0+4S_0}{4}.
 \end{split}
 \label{old-eq:5.3}
\end{equation}
The second numerator is nonnegative.  Since $z_a\le\ell'$, the
energy in this branch is at most
\[
 \begin{aligned}
 &M'+\frac{1+3\ell'-2M'-2\Delta_H+\theta_N}{4}
       +O(\epsilon_0+\kappa_{\mathrm G}+\epsilon_Z)\\
 &\qquad=M'+\frac{d-1/6-2\Delta_H+\theta_N}{4}
       +O(\epsilon_0+\kappa_{\mathrm G}+\epsilon_Z).
 \end{aligned}
\]
After using $O\ge-\epsilon_Z$ and absorbing $\theta_N$, the two
upper bounds are
\[
 \begin{gathered}
 M'-\Delta_H+O(\epsilon_0+\kappa_{\mathrm G}+\epsilon_Z),\\
 M'+\frac{d-1/6-2\Delta_H}{4}
       +O(\epsilon_0+\kappa_{\mathrm G}+\epsilon_Z).
 \end{gathered}
\]
They are nonincreasing in $\Delta_H$.  Since
$\Delta_H\ge-\epsilon_Z$, their formal $\Delta_H=0$ values bound
every actual dyad up to $O(\epsilon_Z)$, whether or not an endpoint
dyad occurs.  Their maximum is therefore at most
$M'+((d-1/6)/4)_++O(\epsilon_0+\kappa_{\mathrm G}+\epsilon_Z)$.
At $d=1/6$ one has $\ell'=z_a=0$ and the clipped unit range falls
in the first branch; bounded negative unit dyads change only
$\epsilon_Z$.  Choose $\epsilon_0$ and $\kappa_{\mathrm G}$ within the prescribed
$\epsilon$ allowance, and then increase the fixed-data lower
threshold so that $\epsilon_Z$ also lies within it.  Summing the
logarithmically many row and dual dyads proves the lemma.
\end{proof}

\subsection{The additive Gram bound}

We next estimate the additive factor in the range $Y'^2/Q\ge1$.  The
full correlation of Lemma~\ref{lem:full-correlation} retains the collision
zeros that determine both its cancellation and its exceptional frequencies.

\begin{proposition}[A quantitative additive Gram bound]\label{prop:probe-gram}
\subsection{Completion of the low bound}

The two preceding estimates now apply to the two factors of the same
rescaled summand in the exact low separation.

\begin{proposition}[The compensated low estimate]\label{prop:probe-low}
For the probe in Equation~\eqref{eq:compensated-probe-definition},
the geometry in Equation~\eqref{old-eq:5.1}, and every
$\epsilon>0$,
\begin{equation}
 |I_{\eta,\mathrm{modified}}(Z)|
       \ll Z^{l_x/2+b/12+\epsilon}=Z^{3/16+\epsilon}.
 \label{old-eq:5.10}
\end{equation}
The exponent is independent of the target and the slot mesh; the
constant and lower threshold may depend on the fixed arithmetic
data, slot system, and smooth tests.
\end{proposition}

\begin{proof}
For a fixed rescaled subset and tuple, put $Q=q_{b_*}X'Y'$.
Its separation is Equation~\eqref{eq:low-separated} with
$A_{m,\sigma,v}(Y')$ and $B^J_{m,\sigma}(Z)$.  Put
$r_J=q_{p_J}/Z^d$; the fixed annular supports give
$0<c_J\le r_J\le C_J<\infty$ uniformly in the tuple.  Exactly,
\[
 \begin{gathered}
 X'=Z^{l_x-d}/r_J,\qquad Y'=Z^{l_y-d}/r_J,\\
 Q=q_{b_*}Z^{M'}/r_J^2,\qquad
 P_a=Y'^2/Q=q_{b_*}^{-1}Z^{1/8}.
 \end{gathered}
\]
In particular $P_a\ge1$ for $Z\ge q_{b_*}^8$.  Enlarge the allowed
fixed-data lower threshold so that $Q,Y'\ge1$ uniformly over the
tuple ratios as well; this is possible because $M'\ge1/2$ and
$l_y-d\ge5/16$.  The remaining length inequalities are
\[
 l_x-d\ge3/16,\qquad l_y-d\ge5/16,
 \qquad l_y-d-11b/6\ge1/12.
\]
The last one gives
\[
 \frac{Y'}{P_a^{11/6}}
 =q_{b_*}^{11/6}r_J^{-1}Z^{l_y-d-11b/6}
       \gg_{\rm fixed}Z^{1/12},
\]
and hence $P_a^2/Y'\ll P_a^{1/6}$.  Proposition~\ref{prop:probe-gram}
therefore gives
\begin{equation}
 \sum_{q_m\ll Q}|A_{m,\sigma,v}(Y')|^2
 \ll (1+|v|)^{J_{\mathrm{Gr}}}(Q/Y')P_a^{1/6}Z^\epsilon.
 \label{old-eq:5.6}
\end{equation}
Cauchy--Schwarz in the row sum, Lemma~\ref{lem:probe-row-norm},
and the integrability of $(1+|v|)^{J_{\mathrm{Gr}}/2}|\widehat W_0(iv)|$ now
bound this tuple's unscaled separation by $Z^\epsilon$ times
\[
 \begin{aligned}
 Q^{-1/2}\bigl[(Q/Y')P_a^{1/6}\bigr]^{1/2}
       \bigl[Z^{M'}\bigr]^{1/2}
 &=\bigl[Z^{M'}/Y'\bigr]^{1/2}P_a^{1/12}\\
 &=r_J(X')^{1/2}P_a^{1/12}.
 \end{aligned}
\]
Since $r_J$ is bounded, this is
$O((X')^{1/2}P_a^{1/12}Z^\epsilon)$.

For $d=\sum_{i\in J}\ell_i$ there are at most $Z^{d+\epsilon}$
rescaled tuples, the coefficient is $O(Z^{-3d/2})$, and
$(X')^{1/2}\asymp X^{1/2}Z^{-d/2}$.  The total additional exponent
is consequently
\begin{equation}
 f(d)=d-3d/2-d/2=-d\le0.
 \label{old-eq:5.5}
\end{equation}
Choose the small powers in the component estimates so that their sum lies
within the prescribed $\epsilon$.  Summing the finitely many subsets proves
Equation~\eqref{old-eq:5.10}.  For the already defined
$C=C_{\mathrm{II}}$ of Equation~\eqref{eq:part-II-mellin-target}, the
normalization used for the principal signal satisfies
\begin{equation}
 \begin{gathered}
 C_{\mathrm{II}}(s)=l_x/2+s-1+h/6=s-11/16,\\
 C_{\mathrm{II}}(7/8)=3/16=l_x/2+b/12.
 \end{gathered}
 \label{old-eq:5.11}
\end{equation}
Thus, under the contradiction assumed in Part~II, this low bound is smaller
than $Z^{C_{\mathrm{II}}(\beta_*)}$ by the exact power
$\Delta=\beta_*-7/8$, apart from the arbitrarily prescribed $\epsilon$.
\end{proof}

\section{The local compensation and its errors}
\label{sec:local-compensation}

The physical modification in Section~\ref{sec:compensation} has already
been estimated from its separated low representation.  We now identify its
full Euler correction and bound the local errors left by the two-term
operation.  Throughout, the local notation $P_p,P_p^*,H_p,V,W,D,R$ is that
of Lemma~\ref{lem:local-euler}; in the shared analytic estimates its variable
$x$ is called $s$.  Fix a nonzero sixth-power-free physical row $u$ with
$(u,S)=1$.  All slot primes belong to the sets $\mathcal P_i(Z)$ of
